
\subsubsection{Model Population}

N=50 models were selected from the Open LLM Leaderboard satisfying the following diversity criteria:

\begin{table}[htbp]
\centering
\resizebox{\columnwidth}{!}{%
\begin{tabular}{ll}
\toprule
Dimension & Coverage \\
\midrule
Architecture & 7 families: Llama, Mistral, Falcon, Phi, Qwen, Gemma, Yi \\
Scale & 4 tiers: 7B, 13B, 34B, 70B \\
Training variant & 4 types: Base, Instruct, Chat, DPO \\
\bottomrule
\end{tabular}
}
\end{table}

This selection ensures correlation estimates are not confounded by population homogeneity. The population excludes proprietary models due to lack of standardized benchmark access.

\subsubsection{Benchmark Score Collection}

Scores were collected for:

\begin{itemize}
\item \textbf{TruthfulQA-MC2:} Multiple-choice accuracy on 817 questions across 38 categories
\item \textbf{HaluEval:} Binary detection accuracy across QA, dialogue, and summarization subtasks (35,000 samples)
\item \textbf{FactScore:} Atomic fact precision on biography generation task
\item \textbf{MMLU:} General knowledge baseline (57 subjects, 14,000 questions)
\end{itemize}

For TruthfulQA and MMLU, evaluation used lm-evaluation-harness with 5-shot prompting. For HaluEval, the official evaluation code was used. For FactScore, proxy methodology validated against available official scores was employed due to computational constraints.

\subsubsection{Correlation Analysis}

Spearman correlation was chosen over Pearson for robustness to non-normal distributions. For each benchmark pair:

\begin{enumerate}
\item Spearman correlation coefficient ($\rho$)
\item 95\% confidence interval via 1,000 bootstrap iterations
\item Statistical significance with Bonferroni correction ($\alpha$ = 0.05/6 = 0.0083)
\end{enumerate}

Correlation interpretation thresholds:
\begin{itemize}
\item r < 0.3: Low correlation (approaching independence)
\item 0.3 $\leq$ r < 0.7: Moderate correlation (partial independence)
\item r $\geq$ 0.7: High correlation (approaching interchangeability)
\end{itemize}

An unrelated-benchmark baseline was established using r(MMLU-Physics, HaluEval) = 0.10.

\subsubsection{Factor Analysis}

Principal Component Analysis was applied with standardized benchmark scores (z-scores). Components with eigenvalue > 1 were extracted. The variance threshold was set at 80\% cumulative variance to determine dimensionality.

\subsubsection{Hypothesis Structure}

Four hypotheses with defined success criteria:

\begin{table}[htbp]
\centering
\resizebox{\columnwidth}{!}{%
\begin{tabular}{lll}
\toprule
Hypothesis & Gate Type & Success Criterion \\
\midrule
H-E1 & MUST\_WORK & 0.10 < r(inter-benchmark) < 0.7 \\
H-M1 & MUST\_WORK & r(TQA, MMLU) < r(MMLU internal) \\
H-M2 & SHOULD\_WORK & r(HE, TQA) < 0.7 \\
H-M3 & SHOULD\_WORK & r(FS, TQA) < 0.7 AND r(FS, HE) < 0.7 \\
\bottomrule
\end{tabular}
}
\end{table}
